% sec-01: the inquiry, the gate, and the funding ladder.  Table 16, Figure 10
% (d2-type layered stack).
\section{The Inquiry, and the Gate in Front of the Pilot}

\subsection{What was asked, and what the answer must be}

An inquiry to apply is an honor, and it is also a question with an exact
answer. The pilot the inquiry concerns is new: NSF announced it in September
2026 as a way to carry deep technologies from its small business portfolio
toward market readiness, with the Institute at the University of Central Florida
selecting the companies and managing their milestones \cite{nsfncti2026}. The
gate to it is published, and the company does not yet pass it. The answer
therefore does three things and no more: it thanks the agency, it states the
gate and the company's place against it in writing, and it asks whether the
inquiry had a different route in mind. The letter to the pilot's principal
investigator asks the same question of the Institute \cite{nctiucf2026}.

\subsection{The tiers and pilots, side by side}

Table 16 sets every NSF tier and pilot the inquiry could concern beside its
amount, its eligibility, and the company's position today. The amounts are
those published in the current solicitations and the pilot's announcement
\cite{nsf26510,nsf26511,nsfncti2026}. Only one row is open to the company now,
and it carries no money: the Project Pitch, which NSF requires before any
Phase I proposal.

\begin{apptable}
\begin{tabularx}{\textwidth}{>{\raggedright\arraybackslash}p{3.9cm} >{\raggedright\arraybackslash}p{4.2cm} >{\raggedright\arraybackslash}p{3.8cm} >{\raggedright\arraybackslash}X}
\headrow \thead{Tier or pilot} & \thead{Amount} & \thead{Who is eligible} & \thead{The company}\\
\midrule
Project Pitch & None; four short fields & Any small business & Submitted today \\
SBIR Phase I & \$305K; 6 to 18 months & Invited after a Pitch & Once invited \\
SBIR Phase II & \$1.25M; about 24 months & Phase I awardees & After Phase I \\
Phase IIB supplement & \$50K to \$500K, with a match & Phase II awardees & Ledger at \$0 \\
NCTI Readiness Pilot & \$1M to \$3.75M each, of \$20M & Recent Phase II or IIB & Not yet eligible \\
Strategic Breakthrough & Up to \$30M; 1 to 1 match & Phase II awardees & Not yet eligible \\
Instrumentation, 26-511 & Share of a \$40M emphasis & Instrument proposers & Outside scope \\
\end{tabularx}
\end{apptable}
\vspace{-0.60cm}
\tabcap{Table 16. The NSF tiers and pilots the inquiry could concern, each with its amount, who\\
is eligible, and the company's place today; only the Project Pitch is open to it now.}

\subsection{The same tiers, drawn as a ladder}

Figure 10 draws Table 16 as the stack it really is. Each layer is reached only
through the one beneath it, and the gate between two layers is an award, a
match, or an invitation. The company stands on the lowest layer. The pilot the
inquiry concerns sits two gates above the first award the company could win.

\begin{appfloat}[!tb]
\begin{appfig}[d2-type / layered stack, each layer gated by the one below]
% Six layers bottom to top, each a D2 container 11.0 cm wide with the tier on
% the left and its amount on the right.  The gate that opens a layer is set in
% the right-hand column at the boundary beneath it.  The company marker is the
% one accent layer, at the bottom.
\tikzset{d2tier/.style={text width=48mm,minimum height=7.2mm,align=left,anchor=west}}
\node[d2key,d2tier,font=\scriptsize\sffamily\bfseries] at (0.00,0.00) {Project Pitch};
\node[d2key,d2tier] at (5.30,0.00) {No funding; four fields; two a year};
\node[d2soft,d2tier,font=\scriptsize\sffamily\bfseries] at (0.00,1.05) {SBIR Phase I};
\node[d2soft,d2tier] at (5.30,1.05) {Up to \$305K, 6 to 18 months};
\node[d2soft,d2tier,font=\scriptsize\sffamily\bfseries] at (0.00,2.10) {SBIR Phase II};
\node[d2soft,d2tier] at (5.30,2.10) {Up to \$1.25M, about 24 months};
\node[d2gray,d2tier,font=\scriptsize\sffamily\bfseries] at (0.00,3.15) {Phase IIB supplement};
\node[d2gray,d2tier] at (5.30,3.15) {\$50K to \$500K, with a match};
\node[d2gray,d2tier,font=\scriptsize\sffamily\bfseries] at (0.00,4.20) {Commercialization Readiness Pilot};
\node[d2gray,d2tier] at (5.30,4.20) {\$1M to \$3.75M each; 6 to 8 firms};
\node[d2gray,d2tier,font=\scriptsize\sffamily\bfseries] at (0.00,5.25) {Strategic Breakthrough};
\node[d2gray,d2tier] at (5.30,5.25) {Up to \$30M, 1 to 1 match};
\begin{scope}[on background layer]
\foreach \y in {0.00,1.05,2.10,3.15,4.20,5.25}{
  \draw[fundmid,line width=0.5pt,rounded corners=3pt,fill=fundwhite] (-0.15,\y-0.47) rectangle (10.75,\y+0.47);}
\end{scope}
\foreach \y/\g in {0.52/invitation after the Pitch, 1.57/a Phase I award, 2.62/a Phase II award, 3.67/Phase II or IIB, 4.72/Phase II and a match}{
  \node[font=\tiny\sffamily,text=fundgrayd,anchor=west] at (11.25,\y) {gate: \g};
  \draw[d2edged] (11.15,\y-0.30) -- (11.15,\y+0.30);}
\node[d2key,text width=24mm,minimum height=7.2mm] (here) at (13.35,0.00) {The company is here};
\draw[d2edgeb] (here.west) -- (10.80,0.00);
\legkey{0.00}{-0.95}{fundaccent}{Open to the company today}
\legkey{4.30}{-0.95}{fundpale}{Reached by winning the layer below}
\legkey{9.05}{-0.95}{fundgrayl}{Needs a Phase II award first}
\ptitle{-0.10}{6.20}{Each layer opens only from the one beneath it}
\end{appfig}
\vspace{-0.60cm}
\figcaption{Figure 10. The NSF tiers and pilots as a layered stack, each opened only by an award\\
or a match on the layer below, with the company on the Pitch layer at the bottom.}
\end{appfloat}

\subsection{The standing fact}

The trial behind the program is built around daraxonrasib, which the FDA
approved on August 26, 2026 as Rasonque for metastatic pancreatic adenocarcinoma
after prior systemic therapy \cite{fdarasonque2026,fdadisco2026,revmedapproval2026}.
The approval followed Revolution Medicines' Phase 1/2 program and its Phase 3
RASolute 302 trial \cite{ctgov001,ctgov302,rasolute302}. ChemicalQDevice's AI
papers on the molecule were developed independently of that program and
simultaneously with it, from the June 2025 identification study
\cite{daraxident} to the Phase 1 protocol and investigational new drug
application of 2026 \cite{onpremwhippl,phase1ind}. The perioperative use the
company's protocol proposes remains investigational, and a separate adjuvant
Phase 3 trial of the molecule has been enrolling since December 2025
\cite{ctgov304}.
